% Part: first-order-logic
% Chapter: natural-deduction
% Section: provability-quantifiers

% pamriksan sipat-sipat provabilitas kang dibutuhake kanggo himpunan
% konsisten maksimal ing bab kelengkapan.

\documentclass[../../../include/open-logic-section]{subfiles}

\begin{document}

\olfileid{fol}{ntd}{qpr}

\olsection{\usetoken{S}{derivability} lan Kuantor}

\begin{explain}
  Teorema kelengkapan uga mbutuhake supaya aturan deduksi natural
  ngasilake kasunyatan ngenani~$\Proves$ kang ditetepake ing bagean
  iki.
\end{explain}

\begin{thm}
\ollabel{thm:strong-generalization} Yen $c$ iku konstanta kang ora
muncul ing $\Gamma$ utawa $!A(x)$ lan $\Gamma \Proves !A(c)$, mula
$\Gamma \Proves \lforall[x][!A(x)]$.
\end{thm}

\begin{proof}
Ayo $\delta$ dadi !!a{derivation} saka $!A(c)$ adhedhasar $\Gamma$.
Kanthi nambahake inferensi \Intro{\lforall}, kita entuk
!!a{derivation} saka $\lforall[x][!A(x)]$. Amarga $c$ ora muncul ing
$\Gamma$ utawa $!A(x)$, syarat eigenvariabel wis dicukupi.
\end{proof}

\begin{prop}
\ollabel{prop:provability-quantifiers}
\begin{tagenumerate}{prvEx,prvAll}
\tagitem{prvEx}{$!A(t) \Proves \lexists[x][!A(x)]$.}{}

\tagitem{prvAll}{$\lforall[x][!A(x)] \Proves !A(t)$.}{}
\end{tagenumerate}
\end{prop}

\begin{proof}
\begin{tagenumerate}{prvEx,prvAll}
\tagitem{prvEx}{Iki !!a{derivation} saka~$\lexists[x][!A(x)]$
  adhedhasar~$!A(t)$:
\begin{prooftree}
\AxiomC{$!A(t)$}
\RightLabel{\Intro{\lexists}}
\UnaryInfC{$\lexists[x][!A(x)]$}
\end{prooftree}}{}

\tagitem{prvAll}{Iki !!a{derivation} saka~$!A(t)$
  adhedhasar~$\lforall[x][!A(x)]$:
\begin{prooftree}
\AxiomC{$\lforall[x][!A(x)]$}
\RightLabel{\Elim{\lforall}}
\UnaryInfC{$!A(t)$}
\end{prooftree}}{}
\end{tagenumerate}
\end{proof}

\end{document}
